% Rubric: Model Specification and Justification.
% Define components, variables, interactions, equations or update rules;
% state assumptions and initialisation; justify model-design and parameter
% choices.

\section{Model specification}
\label{sec:model}

Entity $i$ has state $x_i(t)=(s_i(t),L_i(t))$: open or closed status and carried
load. With neighbours $\mathcal{N}_i$, its generic update is
\begin{equation}
  x_i(t+1) = f\bigl(x_i(t), \{x_j(t)\}_{j \in \mathcal{N}_i}; \theta\bigr),
  \label{eq:update}
\end{equation}
where $\theta$ collects capacities, numerical tolerance and update-mode
parameters. Transport overloads are removed synchronously, then loads are
redistributed or recomputed. The intact-capacity construction follows
\citet{motter2002cascade}; the local variant is specified below. Rounds are
algorithmic, not elapsed time.

\subsection{Components, loads and failure}
The frozen rail layer is an undirected simple tree: 86 stations and 85 adjacent
segments. Multilayer experiments split each station into a
persistent terminal $T{:}i$ and a closable rail facility $R{:}i$ with a
one-minute access edge. Rail edges take
$t_e = \max(0.1,\ d_e/(1000\,v)\cdot 60)$ minutes at an assumed $v = 40$ km/h.
Backup terminal links use four manual effective times (walking, waiting and
transfers) from \texttt{backup\_edges.csv} and GTFS candidates' fastest
scheduled ride times, with no added walking or waiting. Manual rows override
overlapping candidates even when slower, and duplicates within one table keep
the fastest time; the Perth and Perth Underground transfer takes five minutes.
Standby buses activate immediately after the initial closure with unlimited
capacity; multilayer loads are terminal-subset betweenness over the intact
terminal-pair denominator.

The initial load of station $i$ is its unnormalised betweenness on the intact
graph, $L0_i = \sum \sigma_{st}(i) / \sigma_{st}$ over station pairs with
$s \neq i \neq t$, the share of shortest paths through $i$. Unnormalised values
keep rounds comparable as the graph shrinks; normalised loads would rise with
the shrinking node set. The load modes are \texttt{betweenness} as above;
\texttt{betweenness\_freq} with edge length $1/\text{trips}$;
\texttt{betweenness\_plus\_trips} with $0.1\,\text{trips served}_i$ added to
give leaf stations a positive baseline; and \texttt{demand}, the AM-peak
scheduled departure count \texttt{am\_peak\_stops}, a boardings proxy rather
than passenger counts.

Capacities follow
\begin{equation}
  C_i = (1+\alpha) K_i, \qquad \alpha \ge 0,
  \label{eq:capacity}
\end{equation}
with intact reference $K_i = L0_i$ by default, fixed throughout the cascade.
Zero initial load gives zero capacity, but positive static redistribution can
still overload it. For the demand experiment, $L0_i = \text{am\_peak\_stops}_i$,
$f_i = \text{trips\_served}_i$, and $K_i = c f_i$ with
$c = \max_{i:f_i>0}(L0_i/f_i)$. This smallest uniform scale covering demand is
a recorded scenario assumption. Zero frequency with zero demand gives $K_i=0$;
positive demand with zero frequency is rejected.

At each round every surviving station with $L_i > C_i + \delta$ fails, where
the numerical tolerance $\delta$ defaults to $10^{-12}$; the round's overloads
are collected before any removal, and the cascade stops at the first clear
round. The first failure is an explicit \texttt{target} if named, otherwise the
largest initial load (\texttt{load}), the highest degree with station-ID ties
(\texttt{degree}), or a uniform seeded draw (\texttt{random}). A line closure
replaces the trigger with one initial batch, removed before the first check and
excluded from the secondary avalanche size. In the static mode a closed station
forwards its load equally over survivors or in proportion to their capacities,
with equal shares when all have zero; a station with no survivor removes its
load, and otherwise load is conserved. The dynamic mode recomputes loads on the
surviving graph each round and ignores redistribution; on a tree this cannot
raise a survivor's betweenness, so containment reflects dropped journeys rather
than headroom. The containment tolerance $\alpha^*$ is the smallest sampled
$\alpha$ at which the maximum-load trigger causes no secondary failure.

\subsection{Assumptions and initialisation}
All facilities start open; loads and capacities come from the intact graph, and
the trigger or initial batch is applied before the first overload check. The
model assumes undirected adjacent links, scheduled services as load proxies,
hop distances for load concentration, minutes only in the multilayer layer,
symmetric bus times, unlimited immediate standby, and one service date. An
empty initial batch closes nothing; an
$\alpha$ of zero adds no headroom and can fail a station immediately; zero
capacity does not exempt a station from overload checks.

\texttt{src/transperth/config.py} defaults to $\alpha=0.2$, master seed 0,
tolerance $10^{-12}$ and the 81-point grid $\alpha=0,0.025,\ldots,2$.
The tolerance guards floating-point comparisons. The grid starts at zero,
reaches twice the $L0$ reference, and brackets the sampled containment
tolerances near 0.5 and 0.75; its 0.025 step resolves the abrupt static
transition (RQ4).